% P1 学术写作版（2026-11-24）：研究背景与相关工作。
% 由 chapter_02_background_rev_p0_20261124.tex 改写；REV_P0/V1 保持不变。
% 原作者工作出处（原文节名/图号、Supplementary小节、官方仓库文件）详见 REV_P0_EVIDENCE_TABLE A1–A8。未编译。
\section{研究背景与相关工作}
\label{sec:cc-background-p1}

\subsection{空间转录组与空间域识别}
\label{subsec:cc-st-background-p1}
基因表达谱为理解组织细胞状态、发育分化及疾病微环境提供了分子证据，但在不保留空间位置的测序设计中，细胞或组织区域之间的相邻关系难以恢复。空间转录组测序把表达读数与组织中的采样坐标关联起来，使表达相似性和空间结构可以一起分析。空间域识别（spatial domain identification）的目标，是把分子或空间特征相近的观测单元划分为同一区域，所得区域可用于研究组织分层、细胞邻近关系和肿瘤微环境。

空间域聚类与一般的表达聚类既有联系，也有区别。对于表达聚类，两个相距较远的spot只要表达相似，就可以归入同一簇；对于空间域识别，簇在空间上是否连续、簇间交界能否对应真实组织结构，同样具有解释价值。反过来，过分强调空间连续性，也可能让小范围结构和清晰的层间交界变得难以区分。因此，较高的整体聚类一致性并不自动意味着边界定位准确，两者需要分别评价。

人类背外侧前额叶皮层（DLPFC）具有清晰的皮层分层结构，是空间域方法最常用的真实数据基准之一。Maynard等发布了来自3名供体、共12张切片的DLPFC Visium数据，并给出了皮层层与白质的人工参考标注\cite{maynard2021}。使用这一数据时需要注意两点：同一供体的多张切片不是独立的生物学重复；参考标注在层间过渡区域存在不确定性，不应视为完美的绝对真值。

Wu等发布了人类乳腺癌的单细胞和空间转录组图谱，其中包含6例Visium空间样本及病理区域注释\cite{wu2021}。与皮层的层级结构不同，乳腺癌组织区域采用另一套病理注释语义。因此，在两类组织之间比较时，重点应放在同一流程在新队列中的表现和患者异质性上，而不是直接比较两种组织的ARI绝对值。

\subsection{表达降维与空间信息的引入}
\label{subsec:cc-expression-p1}
空间表达矩阵维数高，基因之间的测序深度、零计数比例和技术噪声差异很大，直接在全部基因上聚类既耗费计算，又不稳定。scVI是基于深度生成模型的概率框架，可以从计数数据中学习适合下游分析的低维潜变量\cite{lopez2018}。本研究使用其5维潜变量作为每个spot的基础表示。

把空间上下文引入聚类主要有两类思路。一类先得到表达表示，再聚合空间邻域的特征；另一类直接用空间图或空间统计模型完成联合聚类。SEDR\cite{xu2024}与STAGATE\cite{dong2022stagate}通过图神经网络或自编码器把空间邻接关系纳入表示学习；BANKSY把自身表达与邻域表达加权融合，构造空间增强的特征\cite{singhal2024}；BayesSpace则在低维表达上施加空间先验，用贝叶斯模型完成聚类\cite{zhao2021}。不同方法所用的空间结构（最近邻、物理网格、像素或阵列坐标）各不相同，其原生的数据处理流程也不能互换。因此，比较不同方法时，应在统一的样本和评价掩码下保留各方法自身的输入定义。

\subsection{CellCharter及其在DLPFC上的已有基准}
\label{subsec:cc-cellcharter-p1}
CellCharter由Varrone等提出，是一个用于识别、刻画和比较空间生态位的框架\cite{varrone2024}。设某切片有$N$个spot，每个spot的scVI基础表示为5维，则4层邻域聚合后的特征可写为
\begin{equation}
H=[L_0\mid L_1\mid L_2\mid L_3\mid L_4]\in\mathbb{R}^{N\times25},\qquad L_\ell\in\mathbb{R}^{N\times5},
\label{eq:cc-H-bg-p1}
\end{equation}
其中$L_0$为自身表示，$L_\ell$（$\ell\ge1$）为与该spot相距$\ell$步的邻居的特征均值。CellCharter随后用GMM对拼接后的特征聚类；在簇数未知时，它对每个候选K重复拟合，并比较K、$K-1$、$K+1$三种解之间的Fowlkes--Mallows指数\cite{fowlkes1983}，以聚类稳定性确定K。上述邻域聚合、GMM聚类和稳定性选K都是CellCharter原作者提出的方法，本研究沿用其公开实现。

原作者在同一套12张DLPFC切片上完成了与本研究直接相关的基准工作\cite{varrone2024,cellcharteranalyses2024}：
\begin{itemize}
  \item \textbf{调参与测试。}每名供体随机选1张切片组成调参集，其余9张作为测试集，以ARI评价。
  \item \textbf{超参数网格。}网格包括高变基因数、潜变量维度和邻域层数$\{2,3,4\}$，最终选用5000个高变基因、5维潜变量和4层邻域。
  \item \textbf{性能结论。}单切片聚类中STAGATE表现最好，多样本联合聚类中CellCharter表现最好。
  \item \textbf{选K。}在12张切片的联合数据上，稳定性选K给出9个簇，而BIC的最小值位于19个簇。
\end{itemize}
可见，“邻域层数影响DLPFC聚类ARI”已在原作者的调参中得到体现，只是限于3张调参切片、参考类别数条件和单一ARI指标。至于这一结论在更多切片、不同聚类数规则和不同评价定义下是否依然成立，原作者没有系统考察，这正是本研究的出发点。

多尺度空间信息的构建也并非只有CellCharter一种途径。例如，m2ST研究了多尺度图表示及其融合\cite{zhang2025}。因此，单纯提出“更深的邻域可能有用”或“按尺度拼接特征”，并不足以构成新的算法贡献；更有价值的问题是，尺度效应在什么条件下成立、代价是什么。

\subsection{聚类数选择与成熟基线}
\label{subsec:cc-k-background-p1}
预测簇数K的确定方式会显著影响结果。若直接以人工参考标签的类别数设定K，就使用了标签的类别数先验，本文称之为\textit{Oracle-K}。即使模型拟合时没有读取任何spot的标签，这一设定仍会影响输出，因而不能代表真实的无标签应用。与之相对，\textit{K-blind}指在确定K或分辨率时不使用参考类别信息，例如依据稳定性、BIC等事先规定的规则，或采用固定的分辨率。

不同方法有各自的选K习惯。非空间GMM常用BIC等似然准则，但这类准则容易偏向更多的簇；CellCharter原作者在DLPFC联合数据上也观察到，BIC建议的簇数（19）远多于稳定性峰值（9）\cite{varrone2024}。Leiden是保证社区连通性的图聚类算法\cite{traag2019}，其簇数由分辨率参数间接控制。BayesSpace提供\texttt{qTune}，计算不同候选q的平均伪对数似然，官方建议在\texttt{qPlot}曲线的拐点（elbow）附近选q\cite{bayesspacevignette}。由于这些规则彼此不同，在K-blind条件下比较方法时，必须同时报告各方法实际采用的选K规则。

\subsection{空间聚类的评价指标}
\label{subsec:cc-evaluation-background-p1}
调整兰德指数（ARI）和归一化互信息（NMI）是比较预测聚类与参考标签最常用的指标，两者都不受簇编号的影响。ARI对随机划分作了机会校正，NMI衡量两种划分共享的信息量。但两者衡量的都是观测单元的类别归属，而不是组织边界的几何位置：预测的层间交界略有偏移时，ARI可能几乎不变，而边界检出的结果却会改变。

边界评价可以在Visium六角网格的邻接图上进行：相邻两个spot参考层不同，记为真实边界；预测簇不同，记为预测边界；再由二者的重合计算边界F1。严格匹配要求两者落在同一条网格边上，因此对一两个spot的位置偏移非常敏感；允许1--2步偏移的容差匹配则衡量“邻近交界是否被预测到”。两种定义回答不同的问题，应当并列报告。已有的评价研究也提示，在数据处理、标签质量和模型调节不同时，方法排序具有不确定性\cite{sun2026}。

\subsection{已有研究的不足}
\label{subsec:cc-gap-p1}
综上，CellCharter原作者及多数空间聚类基准的评价，主要建立在Oracle-K条件下的整体一致性指标之上。对于以下问题，现有研究尚未系统回答：
\begin{itemize}
  \item 邻域尺度的收益在不同切片和供体之间是否一致，收益集中在哪些邻域层；
  \item 聚类数规则（参考类别数、固定网格、无标签选K）是否改变尺度结论；
  \item 整体一致性与严格边界定位是否同向变化；
  \item 同一流程在另一组织的患者队列中的适用边界在哪里。
\end{itemize}
本研究针对这些问题开展系统评价。本研究只核对了CellCharter原文、补充材料与官方代码，并未对全领域文献进行系统检索；因此，上述“不足”是相对于CellCharter的已有评价而言的。
